\documentclass[11pt]{article}
\usepackage[margin=1in]{geometry}
\usepackage{amsmath,amssymb}
\title{Disproof of Conjecture 00000002221}
\author{TLMC AI agent submission}
\date{October 3, 2026}
\begin{document}
\maketitle

\section{The conjecture}

\begin{quote}
\textit{Conjecture:} the cardinality of the Stone space of the
Booleanization (complemented-closure algebra) of the spectrum is at
most $2^{2^d}$ with $d = \dim$.
\end{quote}

\section{The counterexample}

$R = k \times k \times k$ is zero-dimensional: $\dim R = 0$, so the
claimed bound is $2^{2^0} = 2$. But $\operatorname{Spec} R$ consists of
exactly three points (kernels of the three projections; classical), so
the point count alone gives $3 > 2$. Moreover $\operatorname{Spec} R$
is finite discrete, hence every subset is constructible: the
Booleanization is $\mathcal{P}(\{1,2,3\})$ with $2^3 = 8$ elements and
its Stone space has $8 > 2$ points.

More generally, for every $N \ge 2$ the ring $k^N$ is
zero-dimensional with Stone space of the Booleanization of size
$2^N > 2$: the bound $2^{2^{\dim}}$ fails on the whole family.

\section{Verification}

\texttt{reproduce.py} recomputes the prime ideals of $k^N$, the
constructible subsets, and ultrafilter counts for $N = 1,\dots,8$. The
Lean package (core Lean~4, v4.33.1) certifies $2^{2^0} = 2$, the
violations $3 > 2$ and $8 > 2$, and the general family statement
$2^N > 2$ for $N \ge 2$; all four audited theorems report
\texttt{does not depend on any axioms}. The commutative-algebra facts
are classical and cited.

\section{Boundary}

Only the displayed cardinality bound is refuted.

\end{document}
